\documentclass[letterpaper,10pt]{article}

%----------PACKAGES----------
\usepackage[
    left=0.50in,
    right=0.50in,
    top=0.48in,
    bottom=0.48in
]{geometry}

\usepackage{hyperref}
\usepackage{titlesec}
\usepackage{xcolor}
\usepackage{enumitem}
\usepackage{tabularx}
\usepackage{array}

%----------COLORS----------
\definecolor{darkblue}{RGB}{0,0,120}

\hypersetup{
    colorlinks=true,
    linkcolor=darkblue,
    urlcolor=darkblue
}

%----------FORMATTING----------
\pagestyle{empty}
\setlength{\parindent}{0pt}
\setlength{\tabcolsep}{0pt}
\renewcommand{\baselinestretch}{1.00}

\titleformat{\section}
  {\large\bfseries\uppercase}
  {}
  {0em}
  {}
  [\titlerule]

\titlespacing*{\section}{0pt}{3pt}{2pt}

%----------CUSTOM COMMANDS----------

\newcommand{\cvItem}[1]{
    \item \small{#1}
}

\newcommand{\cvSubheading}[4]{
    \vspace{1pt}
    \begin{tabularx}{\textwidth}{X r}
        \textbf{#1} & #2 \\
        \textit{\small #3} & \small #4 \\
    \end{tabularx}
    \vspace{-3pt}
}

\newcommand{\cvProjectHeading}[2]{
    \vspace{1pt}
    \begin{tabularx}{\textwidth}{X r}
        \textbf{#1} & \small #2 \\
    \end{tabularx}
    \vspace{-3pt}
}

\newcommand{\cvItemListStart}{
    \begin{itemize}[
        leftmargin=0.16in,
        itemsep=0.75pt,
        parsep=0pt,
        topsep=1.5pt,
        partopsep=0pt
    ]
}

\newcommand{\cvItemListEnd}{
    \end{itemize}
    \vspace{0pt}
}

\begin{document}

%===========================================================
% HEADER
%===========================================================

\begin{center}
    {\Huge \textbf{Tushar Nayak}} \\[-1pt]

    \small
    \href{mailto:tusharnayak@outlook.com}{tusharnayak@outlook.com}
    $|$
    \href{tel:+14124785677}{(412) 478-5677}
    $|$
    \href{https://linkedin.com/in/nayaktushar}{linkedin.com/in/nayaktushar}
    $|$
    \href{https://tushar-nayak.github.io}{tushar-nayak.github.io}
    $|$
    \href{https://github.com/tushar-nayak}{github.com/tushar-nayak}
\end{center}

\vspace{-6pt}


%===========================================================
% RESEARCH FOCUS
%===========================================================

\section{Research Focus}

\small{
Biomedical engineering researcher developing geometry-aware learning methods
for medical image analysis, with an emphasis on recovering 3D anatomy, motion,
and deformation from sparse, noisy, or multimodal clinical observations.
Research interests include sparse-view 2D--3D reconstruction, deformable
anatomical modeling and registration, neural implicit and Gaussian
representations, graph-based geometric learning, and physics- and
geometry-informed methods for image-guided intervention.
}


%===========================================================
% EDUCATION
%===========================================================

\section{Education}

\cvSubheading
    {Carnegie Mellon University}
    {Pittsburgh, PA}
    {M.S. Biomedical Engineering, Computational Research Track}
    {May 2026}

\small{
\textit{Courses:
Computer Vision, Learning for 3D Vision, Medical Image Analysis,
Image-Based Computational Modeling and Analysis, Clinical Translations of AI.
}}

\vspace{2pt}

\cvSubheading
    {Manipal Institute of Technology}
    {Manipal, India}
    {B.Tech. Biomedical Engineering, Minor in Data Science}
    {May 2023}


%===========================================================
% RESEARCH EXPERIENCE
%===========================================================

\section{Research Experience}


%---------------- UCHICAGO ----------------

\cvProjectHeading
    {Staff Scientist, Surgical Biomechanics Lab, University of Chicago}
    {Aug 2026 -- Present}

\textit{\small
Geometry-Aware Learning and Computational Modeling for Vascular Biomechanics
}

\cvItemListStart

    \cvItem{
    Serving as staff scientist and computer vision/ML researcher under
    Dr. Luka Pocivavsek, developing computational imaging and machine learning
    methods for patient-specific vascular anatomy, biomechanics, and
    image-guided intervention.
    }

    \cvItem{
    Building CT/CTA-to-3D vascular modeling pipelines spanning medical-image
    preprocessing, segmentation, surface reconstruction, mesh processing,
    anatomical correspondence, and geometry quality assurance.
    }

    \cvItem{
    Investigating longitudinal vascular remodeling through regional geometric
    descriptors, including baseline integrated Gaussian curvature and
    subsequent patch-wise surface-area growth, with extension toward
    biomechanical stress analysis.
    }

\cvItemListEnd


%---------------- PITT / TRACER ----------------

\cvProjectHeading
    {Research Collaborator, Surreality Lab, University of Pittsburgh}
    {Sept 2026 -- Present}

\textit{\small
Multimodal AI for Intraoperative Sensing and Image-Guided Intervention
}

\cvItemListStart

    \cvItem{
    Contributing to \textsc{TRACER}, an automated intraoperative speech-analysis
    pipeline for awake craniotomy that isolates patient speech from noisy
    operating-room audio using speaker attribution, word-onset localization,
    overlap detection, and neural speech separation.
    }

    \cvItem{
    Developing and evaluating deep learning methods for robust speech-overlap
    detection and selective source separation, with emphasis on identifying
    failure modes and preventing enhancement models from degrading clinically
    relevant speech measurements.
    }

    \cvItem{
    Auditing robustness and deployment assumptions across diarization,
    temporal localization, overlap detection, and speech-separation components
    to improve reliability for quantitative intraoperative language assessment.
    }

\cvItemListEnd


%---------------- FLUOROGRID ----------------

\cvProjectHeading
    {Graduate Researcher, Computational Engineering \& Robotics Lab, Carnegie Mellon University}
    {Aug 2024 -- Jul 2026}

\textit{\small
Real-Time 3D Vascular Deformation Estimation from Sparse 2D Angiography
}

\cvItemListStart

    \cvItem{
    Developed \textsc{FluoroGRID}, a real-time multimodal framework for
    estimating patient-specific 3D vascular deformation from monocular
    fluoroscopy by combining image evidence, preoperative anatomy, camera
    geometry, and sparse vascular control points.
    }

    \cvItem{
    Architected a quad-encoder geometric deep learning system integrating
    ResNet image features, ordered-chain graph representations, sampled
    vascular mesh geometry, and fluoroscopic acquisition metadata for
    physically plausible 3D deformation estimation.
    }

    \cvItem{
    Designed geometry-aware objectives incorporating differentiable projection,
    structural smoothness, curvature constraints, and discrete differential
    geometry; achieved sub-millimeter reconstruction precision with
    real-time inference.
    }

    \cvItem{
    Built a physics-grounded synthetic data pipeline for vascular deformation
    and fluoroscopy-style projection, generating more than 370,000 paired
    3D--2D training examples across diverse viewing geometries and
    deformation conditions.
    }

\cvItemListEnd


%---------------- GBM ----------------

\cvProjectHeading
    {Research Assistant, The $\forall$ Lab \& Image Science Lab, Carnegie Mellon University}
    {Feb 2025 -- Aug 2025}

\textit{\small
Longitudinal Glioblastoma Evolution and Visual Forecasting
}

\cvItemListStart

    \cvItem{
    Developed a continuous-time longitudinal forecasting framework combining
    Attention U-Net representations, temporal embeddings, and Neural ODE
    dynamics to predict future glioblastoma state from multimodal MRI.
    }

    \cvItem{
    Reduced mean-squared prediction error by 6.7\% relative to a persistence
    baseline and reduced relative FLAIR-volume error by 98\%; presented the
    work at the 2025 CMU Biomedical Engineering Forum.
    }

\cvItemListEnd


%===========================================================
% SELECTED RESEARCH PROJECTS
%===========================================================

\section{Selected Research Projects}


%---------------- CARDIAC ----------------

\cvProjectHeading
    {\href{https://github.com/tushar-nayak/cardiac-reconstruction-evolved}
    {Sparse Echocardiography to 3D Cardiac Volume Reconstruction}}
    {Aug 2025 -- Present}

\cvItemListStart

    \cvItem{
    Developed a coordinate-based neural implicit reconstruction framework that
    recovers dense 3D left-ventricular occupancy from only three sparse 2D
    echocardiography slices, achieving 0.864 3D Dice and 0.765 3D IoU on
    \textsc{MITEA-CAP}.
    }

    \cvItem{
    Integrated Fourier positional encoding, differentiable multi-view
    projection, learnable \textsc{SE(3)} slice-pose refinement, volumetric
    regularization, and Reptile meta-initialization for subject-specific
    sparse-view reconstruction.
    }

    \cvItem{
    Extended the framework toward explicit Gaussian occupancy representations
    with non-saturating primitive aggregation, reaching 0.950 IoU and
    0.966 sampled occupancy accuracy in preliminary experiments.
    }

\cvItemListEnd


%---------------- ENDOSCOPIC 3DGS ----------------

\cvProjectHeading
    {\href{https://github.com/tushar-nayak/endo-splat}
    {Deformable 3D Reconstruction of Live Endoscopic Surgical Scenes}}
    {Apr 2026 -- Present}

\cvItemListStart

    \cvItem{
    Prototyping deformable 3D Gaussian Splatting for dynamic monocular
    endoscopic reconstruction using depth/pose-based initialization,
    differentiable rendering, deformation fields, and camera-aware projection.
    }

    \cvItem{
    Exploring semantic feature alignment for open-vocabulary querying and
    interpretation of reconstructed surgical scenes.
    }

\cvItemListEnd


%---------------- MAMMOGRAPHY ----------------

\cvProjectHeading
    {\href{https://github.com/tushar-nayak/grading-cbisddsm}
    {Geometric Multi-View Breast Cancer Detection}}
    {Jan 2026 -- May 2026}

\cvItemListStart

    \cvItem{
    Developed a geometric alignment and dual-view correspondence framework for
    paired CC/MLO mammography, integrating lesion segmentation, localization,
    image registration, and cross-view attention.
    }

    \cvItem{
    Reduced cross-view lesion-centroid discrepancy by 99.1\%, enabling
    spatially aligned multi-view representations for downstream malignancy
    classification.
    }

\cvItemListEnd


%===========================================================
% PUBLICATIONS
%===========================================================

\section{Selected Publications, Manuscripts \& Presentations}

\cvItemListStart

    \cvItem{
    \textbf{Nayak, T.}, Basdeo, R., Andrews, E., et al.
    \textit{Real-time 3D Neuro-Vascular Deformation Estimation from Sparse
    2D Angiography}. Manuscript in preparation.
    }

    \cvItem{
    \textbf{Nayak, T.}
    \textit{Approaches in Sparse-View 2D Echocardiography to 3D Cardiac
    Volume Reconstruction using Neural Implicit Representations and Gaussian
    Occupancy Fields}. Manuscript in preparation.
    }

    \cvItem{
    Basdeo, R., Patel, S., \textbf{Nayak, T.}, et al.
    \textit{A Review of Robotics in Achieving Autonomous and Remote
    Neuroendovascular Intervention}. Under review, Journal of Intelligent
    \& Robotic Systems.
    }

    \cvItem{
    Kite, T., \textbf{Nayak, T.}, Jaffee, S., et al. (2026).
    \textit{The use of artificial intelligence and machine learning to predict
    tumor recurrence in high-grade gliomas: a systematic review}.
    Neurosurgical Review, 49(1), 489.
    }

    \cvItem{
    \textbf{Nayak, T.}, Gokulkrishnan, N., Chadaga, K., Sampathila, N.,
    et al. (2024).
    \textit{Automated histopathological detection and classification of lung
    cancer with an image pre-processing pipeline and spatial attention with
    deep neural networks}. Cogent Engineering, 11, 2357182.
    }

    \cvItem{
    \textbf{Nayak, T.}, Sampathila, N., \& Gokulkrishnan, N. (2023).
    \textit{Processing and Detection of Lung and Colon Cancer from
    Histopathological Images using Deep Residual Networks}.
    IEEE CONECCT 2023.
    }

    \cvItem{
    \textbf{Nayak, T.}, Basdeo, R., Andrews, E., et al. (2026).
    \textit{Toward Real-Time 3D Vascular State Estimation from Monocular
    Fluoroscopy for Endovascular Robotic Neurosurgery}.
    2nd International Workshop on the Evolving Landscape of Surgical Robotics,
    IEEE/RSJ IROS 2026.
    }

\cvItemListEnd


%===========================================================
% TECHNICAL SKILLS
%===========================================================

\section{Technical Skills}

\small{

\textbf{Programming \& Scientific Computing:}
Python, C++, MATLAB, R, Julia, CUDA, Linux, Git, Docker, Bash, \LaTeX. 

\textbf{Deep Learning:}
PyTorch, PyTorch3D, PyTorch Geometric, MONAI, TorchDiffEq, TensorFlow/Keras;
GNNs, U-Nets, Transformers, Neural ODEs, PINNs, neural implicit
representations, Gaussian representations. 

\textbf{Medical Imaging \& Computer Vision:}
CT, MRI, fluoroscopy/angiography, ultrasound/echocardiography, endoscopy;
segmentation, registration, 2D--3D reconstruction, deformable modeling,
multi-view geometry, camera geometry, differentiable projection/rendering.

\textbf{3D Geometry \& Imaging Software:}
OpenCV, SimpleITK/ITK, VTK, VMTK, nnU-Net;
point clouds, surface meshes, anatomical geometry processing, DICOM, NIfTI.

}


%===========================================================
% TEACHING & HONORS
%===========================================================

\section{Teaching \& Honors}

\small{

\textbf{Graduate Teaching Assistant, Carnegie Mellon University}
\hfill 2025--2026

Applied Deep Learning; Computer Vision for Engineers;
Machine Learning Applications in Experimental Biomedical Engineering Research;
Fundamentals of Computational Biomedical Engineering.

\vspace{3pt}

\textbf{Honors:}
Carnegie Mellon University Biomedical Engineering Department Head's Fellowship
(2024--2026);
Best Paper Award, 2nd International Conference on Artificial Intelligence,
Computational Electronics and Communication (2023).

}

\end{document}